El presente trabajo de graduación evalúa el impacto de la inteligencia artificial generativa en el desarrollo de videojuegos, analizando la diferencia en la experiencia del jugador y la percepción de calidad entre el código programado de forma tradicional y el código asistido por IA en Godot Engine. El estudio aborda la problemática de determinar si la integración de modelos generativos en la ingeniería de software degrada la calidad funcional del producto o si las variaciones en su evaluación responden a sesgos cognitivos sobre la autoría de la tecnología.

Para ello, se implementó un diseño experimental A/B con una muestra multidisciplinaria de 60 estudiantes universitarios, quienes evaluaron prototipos bidimensionales mediante el cuestionario estandarizado Player Experience Inventory (PXI) bajo condiciones ciegas e informadas. Los resultados cuantitativos demostraron paridad funcional y técnica en la prueba a ciegas, no registrándose diferencias significativas en la experiencia de juego. No obstante, la revelación explícita de la autoría para el grupo informado provocó un decremento en las calificaciones hacia la versión hecha con IA.

Se concluye que la experiencia creada por IA alcanza un rendimiento equivalente al desarrollo humano, pero su percepción se encuentra significativamente condicionada por un sesgo de etiquetado y por las actitudes previas del usuario, lo que permitió aceptar la hipótesis de investigación. Se recomienda extender el estudio a audiencias especializadas e industrias comerciales, implementar evaluaciones longitudinales e incorporar condiciones de prueba con información parcial para discriminar la capacidad real de detección frente al sesgo inducido. 